\begin{theorem}[Wild odd scalar phase products]\label{mod:cases-wildodd-scalarphaseproducts}
Let $p\ne2$ be prime, let $k$ be a finite field of characteristic
$p$, let $\psi:k\to\mathbf C^\times$ be a nontrivial additive
character, and write
\[
 \nu=\nu_k,\qquad q_0=Q_\psi(1,0).
\]
The basic phase has the three exact powers
\[
 q_0^2=\nu(-1),\qquad
 q_0^{p-1}=\nu(-1),\qquad
 q_0^p=q_0\nu(-1).
\]

If $a,c\in k^\times$ and $b\in k$, simultaneous scalar change of the
two coefficients gives
\[
 Q_\psi(ca,cb)
 =\nu(c)\psi\!\left(-\frac{(c-1)b^2}{2a}\right)Q_\psi(a,b).
\]
More generally, for a finite index set $I$, coefficients
$a_i\in k^\times$, and $b_i\in k$, one has
\[
 \prod_{i\in I}Q_\psi(a_i,b_i)
 =\nu\!\left(\prod_{i\in I}a_i\right)
  \psi\!\left(\sum_{i\in I}-\frac{b_i^2}{2a_i}\right)
  q_0^{|I|}.
\]
Thus every displayed denominator is protected by the explicit condition
$a_i\ne0$.

For $a,d\in k^\times$ and $\gamma\in k$, the repeated phase and the
upper homogeneous scalar phase are
\[
 Q_\psi(a,a\gamma)^p=q_0\nu(-a),
 \qquad
 Q_\psi(-a d^{p-1},0)=q_0\nu(-a),
\]
and hence, with no cancellation of the quadratic-character factor,
\[
 Q_\psi(a,a\gamma)^p=Q_\psi(-a d^{p-1},0).
\]
Equivalently, with the zero index included, the repeated-phase product
over the whole prime field is
\[
 \prod_{j\in\F_p}Q_\psi(a,a\gamma)
 =Q_\psi(-a d^{p-1},0).
\]
Here $a\ne0$ is the quadratic denominator condition and $d\ne0$
ensures that the upper quadratic coefficient is nonzero.  The first
identity uses that every value of $\psi$ has $p$-th power one; the
second uses the even exponent $p-1$ only to prove
$\nu(d^{p-1})=1$, while retaining $\nu(-1)$.

Finally let
$\omega:\F_p^\times\to k^\times$ be an injective homomorphism.  For
every $a\in k^\times$ and $\gamma\in k$,
\[
 \prod_{j\in\F_p^\times}
 Q_\psi\bigl(a\omega(j),a\omega(j)\gamma\bigr)=1.
\]
The index is $\F_p^\times$, so the zero scalar is omitted.  Indeed,
\[
 \sum_{j\in\F_p^\times}\omega(j)=0,
 \qquad
 \prod_{j\in\F_p^\times}\omega(j)=-1.
\]
The additive phase therefore has argument zero, whereas the two
remaining scalar factors are separately
\[
 \nu\!\left(a^{p-1}(-1)\right)=\nu(-1),
 \qquad q_0^{p-1}=\nu(-1).
\]
Their product is $\nu(-1)^2=1$.  This is only the nonboundary scalar
calculation; no affine-pencil boundary identity is asserted or used.
\LeanFile{LanglandsFirstMainLemma/Cases/WildOdd/ScalarPhaseProducts.lean}
\lean{LanglandsFirstMainLemma.wildOdd_basicPhasePowers}
\lean{LanglandsFirstMainLemma.wildOdd_quadraticPhase_scale}
\lean{LanglandsFirstMainLemma.wildOdd_quadraticPhase_product}
\lean{LanglandsFirstMainLemma.wildOdd_repeatedPhase}
\lean{LanglandsFirstMainLemma.wildOdd_upperScalarPhase}
\lean{LanglandsFirstMainLemma.wildOdd_repeatedPhase_eq_upperScalar}
\lean{LanglandsFirstMainLemma.wildOdd_repeatedPhase_primeFieldProduct}
\lean{LanglandsFirstMainLemma.wildOdd_scalarPencil}
\leanok
\SourceText{Lemma lem:q0-powers; Lemma lem:quadratic-phase-new; Proposition prop:odd-phase-cancel.}
\uses{mod:finitefield-quadraticphase,mod:finitefield-q0powers,mod:basic-finiteproducts}
\FormalizationContract
\end{theorem}
\begin{proof}
The three powers of $q_0$ are Lemma lem:q0-powers in the explicit
characteristic-$p$ notation.  The scalar and finite-product identities
follow by substituting Lemma lem:quadratic-phase-new and multiplying the
quadratic-character and additive-character factors without changing
their order or sign.

For the repeated phase, write
$z=-(a\gamma)^2/(2a)$.  Characteristic $p$ gives
$\psi(z)^p=\psi(pz)=1$.  Since $p$ is odd,
$\nu(a)^p=\nu(a)$, while $q_0^p=q_0\nu(-1)$, so
\[
 Q_\psi(a,a\gamma)^p
 =\nu(a)q_0\nu(-1)=q_0\nu(-a).
\]
Also $p-1$ is even, hence $\nu(d^{p-1})=1$; the nonzero coefficient
$-ad^{p-1}$ consequently gives the stated upper phase.

For the prime-field product, the finite-product formula leaves the
coefficient product $a^{p-1}(-1)$, zero additive argument, and the
power $q_0^{p-1}$.  Since $p-1$ is even,
$\nu(a)^{p-1}=1$.  The coefficient character is therefore exactly
$\nu(-1)$, not $1$; multiplying it by
$q_0^{p-1}=\nu(-1)$ gives one.  No boundary formula enters the proof.
\end{proof}
